
\section{Ordination Refusal and the Constitution of the United Church}

Editorial article in "Folkebladet" for January 16, 1895.

The Constitution of the United Church is, both in form and in content, a contract or agreement between Lutheran congregations. It is not intended to set an authority over the congregations; rather, it binds several congregations together for certain purposes, whereby those congregations thus receive greater power and ability to carry out those tasks which seem too burdensome for a single congregation. Among such tasks are named the education of pastors, home mission, the distribution of the Bible, and mission to the heathen.

The ordination of pastors also is in our Constitution made a matter of the fellowship, and of course for the same reason as the other matters that belong among the fellowship’s tasks. It is only because it is expedient, not because it is necessary. That is the whole reason and the only reason. It is not assigned to the fellowship because the congregation does not have the right to do it itself; but because there are certain advantages, both for the congregations and for the pastors, in several congregations working together in that matter.

That something is set forth in the Constitution as a matter of the fellowship does not make it a sin for a congregation to concern itself with it. It is not even certain that it becomes disorder if the congregation concerns itself with it. An example will at once show this: home mission is a matter of the fellowship, and the Constitution treats it thus when it counts among the purposes of the fellowship "to gather and organize Lutheran congregations." In that there lies no prohibition against a single congregation with its pastor carrying on home mission and "gathering and organizing" another Lutheran congregation in its immediate vicinity or farther away. This example shows altogether clearly that it is by no means the intention that the work of the fellowship should exclude, hinder, and forbid the work of the congregations, each by itself; that a matter is set forth as a matter of the fellowship does not hinder a congregation from concerning itself with the same matter. On the contrary, the fellowship can do only what it does because the congregation has the right and duty according to the Word of God to do it.

So also with ordination. It is an act which each congregation by itself can perform with its chosen pastor. No fellowship is needed for that matter’s sake. When, then, our Constitution names ordination as a matter of the fellowship, that is only for expediency’s sake. For when several congregations cooperate in this matter, it is thereby said that the man whom the one congregation has chosen as its pastor is acknowledged as a Lutheran pastor also by several other congregations. He will by them all be esteemed as a servant of the Lord, and the examination which fittingly goes before ordination therefore need not be repeated if another congregation calls him as pastor, just as little as the ordination itself need be repeated.

This is a convenient and reasonable way of ordering things; but it does not alter the matter itself: ordination is a congregational act, which is carried out by several congregations in common only for reasons of expediency.

From this it follows with complete necessity that ordination is not such a matter in which the fellowship is lord over the congregations and their right of call; rather, it is only a matter in which the fellowship serves the congregations.

Our Constitution is far from perfect; yet with good will one could well make use of it without crushing the little sprout of congregational freedom, which as it were peeped forth in the time of union. But when a number of partisan pastors wish to misuse the Constitution in order to abolish congregational freedom, it will come to pass that it is found true that he who sows wind reaps storm, as is just and proper.

Georg Sverdrup
